\documentclass[a4paper,11pt]{jsarticle}
\usepackage[dvipdfmx]{graphicx}
\usepackage{geometry}
\usepackage{amsmath}
\usepackage{amssymb}
\usepackage{booktabs}
\geometry{margin=22mm}

\title{p2MEG における停止ミューオン数正規化の導出}
\author{p2MEG 解析メモ}
\date{2026年3月6日}

\begin{document}
\maketitle

\section{独立な記号}
\begin{table}[h]
\centering
\begin{tabular}{ll}
\toprule
記号 & 定義 \\
\midrule
$j$ & run または幾何配置の添字 \\
$N_{j,\mathrm{HE}}^{\mathrm{data}}$ & 実験で得た高エネルギー窓の $E_e$ 1次元カウント \\
$\kappa_{j}^{\mathrm{Michel}}$ &
\(
\kappa_{j}^{\mathrm{Michel}}
=
\dfrac{N_{j,\mathrm{HE}}^{\mathrm{Michel,MC}}}{N_{\mu,j}^{\mathrm{stop,MC}}}
\) \\
$\kappa_{j}^{\mathrm{sig}}$ &
\(
\kappa_{j}^{\mathrm{sig}}
=
\dfrac{N_{j}^{\mathrm{sig,MC}}}{N_{\mu,j}^{\mathrm{stop,MC}}}
\) \\
\bottomrule
\end{tabular}
\end{table}

\section{導出}
\subsection{有効停止ミューオン数の定義}
signal 事象数を $N_{j}^{\mathrm{sig}}$、分岐比を $\mathcal{B}\equiv \mathrm{BR}(\mu\to e\gamma)$ とする。
run $j$ ごとに
\begin{equation}
N_{j}^{\mathrm{sig}}
=
N_{\mu,j}^{\mathrm{eff}}\,\mathcal{B}
\end{equation}
となるように $N_{\mu,j}^{\mathrm{eff}}$ を定義する。

従って、全 run を合わせると
\begin{equation}
N^{\mathrm{sig}}
=
\sum_j N_{j}^{\mathrm{sig}}
=
\mathcal{B}\sum_j N_{\mu,j}^{\mathrm{eff}}
\end{equation}
であり、
\begin{equation}
\mathcal{B}
=
\frac{N^{\mathrm{sig}}}{N_{\mu}^{\mathrm{eff}}},
\qquad
N_{\mu}^{\mathrm{eff}} \equiv \sum_j N_{\mu,j}^{\mathrm{eff}}
\label{eq:br_def}
\end{equation}
となる。

\subsection{高エネルギー Michel データによる導出}
高エネルギー窓の実験カウント $N_{j,\mathrm{HE}}^{\mathrm{data}}$ は、停止ミューオン数を $N_{\mu,j}^{\mathrm{stop}}$、
$e$ が検出器で反応して再構成される確率を $\eta_{e,j}$ とすると、
\begin{equation}
N_{j,\mathrm{HE}}^{\mathrm{data}}
=
N_{\mu,j}^{\mathrm{stop}}
\times
\left(
\frac{N_{j,\mathrm{HE}}^{\mathrm{Michel,MC}}}{N_{\mu,j}^{\mathrm{stop,MC}}}
\right)
\times
\eta_{e,j}
\end{equation}
と書ける。

表の定義
\begin{equation}
\kappa_{j}^{\mathrm{Michel}}
=
\frac{N_{j,\mathrm{HE}}^{\mathrm{Michel,MC}}}{N_{\mu,j}^{\mathrm{stop,MC}}}
\end{equation}
を用いると、
\begin{equation}
N_{j,\mathrm{HE}}^{\mathrm{data}}
=
N_{\mu,j}^{\mathrm{stop}}
\times
\kappa_{j}^{\mathrm{Michel}}
\times
\eta_{e,j}
\end{equation}
となる。従って、
\begin{equation}
N_{\mu,j}^{\mathrm{stop}}
=
N_{j,\mathrm{HE}}^{\mathrm{data}}
\times
\frac{1}{\kappa_{j}^{\mathrm{Michel}}}
\times
\frac{1}{\eta_{e,j}}
\label{eq:nstop}
\end{equation}
を得る。

次に、$\gamma$ が検出器で反応して再構成される確率を $\eta_{\gamma,j}$ とすると、
\begin{equation}
N_{\mu,j}^{\mathrm{eff}}
=
N_{\mu,j}^{\mathrm{stop}}
\times
\left(
\frac{N_{j}^{\mathrm{sig,MC}}}{N_{\mu,j}^{\mathrm{stop,MC}}}
\right)
\times
\eta_{e,j}
\times
\eta_{\gamma,j}
\end{equation}
である。

表の定義
\begin{equation}
\kappa_{j}^{\mathrm{sig}}
=
\frac{N_{j}^{\mathrm{sig,MC}}}{N_{\mu,j}^{\mathrm{stop,MC}}}
\end{equation}
を用いると、
\begin{equation}
N_{\mu,j}^{\mathrm{eff}}
=
N_{\mu,j}^{\mathrm{stop}}
\times
\kappa_{j}^{\mathrm{sig}}
\times
\eta_{e,j}
\times
\eta_{\gamma,j}
\label{eq:neff0}
\end{equation}
となる。

ここで式\eqref{eq:nstop} を式\eqref{eq:neff0} に代入すると、
\begin{equation}
N_{\mu,j}^{\mathrm{eff}}
=
N_{j,\mathrm{HE}}^{\mathrm{data}}
\times
\frac{1}{\kappa_{j}^{\mathrm{Michel}}}
\times
\frac{1}{\eta_{e,j}}
\times
\kappa_{j}^{\mathrm{sig}}
\times
\eta_{e,j}
\times
\eta_{\gamma,j}
\end{equation}
となり、$e$ 側の確率 $\eta_{e,j}$ は打ち消し合って
\begin{equation}
N_{\mu,j}^{\mathrm{eff}}
=
N_{j,\mathrm{HE}}^{\mathrm{data}}
\times
\frac{\kappa_{j}^{\mathrm{sig}}}{\kappa_{j}^{\mathrm{Michel}}}
\times
\eta_{\gamma,j}
\label{eq:neff_final}
\end{equation}
を得る。

最後に、全 run を合わせた有効停止ミューオン数は
\begin{equation}
N_{\mu}^{\mathrm{eff}}
=
\sum_j N_{\mu,j}^{\mathrm{eff}}
=
\sum_j
N_{j,\mathrm{HE}}^{\mathrm{data}}
\times
\frac{\kappa_{j}^{\mathrm{sig}}}{\kappa_{j}^{\mathrm{Michel}}}
\times
\eta_{\gamma,j}
\end{equation}
である。

\section{結果}
従って、Michel 専用 MC を用いて高エネルギー $E_e$ 1次元データで正規化すると、
\begin{equation}
\boxed{
N_{\mu,j}^{\mathrm{eff}}
=
N_{j,\mathrm{HE}}^{\mathrm{data}}
\times
\frac{\kappa_{j}^{\mathrm{sig}}}{\kappa_{j}^{\mathrm{Michel}}}
\times
\eta_{\gamma,j}
}
\end{equation}
となる。

したがって、分岐比は式\eqref{eq:br_def} により
\begin{equation}
\mathrm{BR}(\mu\to e\gamma)
=
\frac{N^{\mathrm{sig}}}{N_{\mu}^{\mathrm{eff}}}
\end{equation}
で与えられる。

\end{document}
